% Modernised reading — fonds Alexandre Grothendieck, shelfmark 104, pages 1-16.
%
% This is an INTERPRETATION, not a transcription. Everything the critical
% apparatus carried moves into footnotes. In any dispute of substance the
% transcription governs: whoever cites Grothendieck must cite batch-01.fr.tex.
%
% Pass: Opus 5 (claude-opus-5), 2026-09-13 — first pass, unchecked against the
% pages by a human. The folder was read whole; it holds one batch, pages 1-16,
% all transcribed.
%
% What the folder is. A category M of « modèles » in which every arrow is a
% monomorphism and every endomorphism an identity is replaced by an
% order-theoretic datum — the poset N of isomorphism classes, the posets I_n
% of subobjects, the type maps tau_n and the admissible maps u_xi — from which
% M is recovered up to equivalence (pages 6-10); three examples, which are
% the simplex, the cube and the globe (page 10); a question on regular
% spherical posets (page 11); and the presheaves on M that are unions of
% their model subobjects, described by « M-pinned » posets (pages 12-16).
% Pages 1-5 are loose sheets bearing on the same construction and on the
% globular apparatus of Pursuing Stacks, read briefly; nothing dates them
% relative to pages 6-16.
%
% Notation fixed for the whole document. M the category of models, M^ its
% presheaves (he writes C^ on pages 14 and 16 for the same thing). N (his
% barred N) the poset of isomorphism classes, distinct from the natural
% numbers (his plain N). D_n a model of class n, I_n = D~_n its poset of
% subobjects, d_n its greatest element, tau_n : I_n -> N the type map,
% I_{n,n'} = tau_n^{-1}(n'), u^n_xi : I_{n'} -> I_n the admissible map of xi.
% For an admissible presheaf X, K = X~ the poset of its model subobjects,
% tau_K and u^K_xi its pinning.
%
% Substantive departures from the pages, each footnoted where it occurs:
% — page 9: « rigide (deux objets isom. sont identiques) » is what is now
%   called skeletal; the reading uses that word and says so.
% — page 9: the list of data omits u_{d_n} = id and tau_n o u_xi = tau_{n'};
%   both follow from the composition law, and the reading says how.
% — page 10, Ex. 3: I_n = (Delta_n x {±1}) ⊔ {d_n} has 2n+3 elements; the
%   margin's 2n+1, and the geometry of the hemispherical disc, want
%   Delta_{n-1}. And « {±1} muni de l'ordre chaotique » gives a preorder; the
%   order meant is (k, e) < (k', e') iff k < k'. The reading states both.
% — page 13: the triangle's right vertex is L_{<=y}, not K_{<=y}.
% — page 5: « pi_0(F_*, x) = classes de lacets en x » names a pi_1.
% — pages 2-3: F_0 = Ob Delta_n (not Delta_{n'}), third line Fl C_1, last
%   factor Ob C_{n+1} (struck on the page), as the pattern of the formulas
%   requires; the amalgamation over F-dot x I (the transcription reads a
%   B-dot, flagged uncertain).
% — page 4: the page's leftmost term Z[F_{n-1}], repeated, is dropped.
% — page 11, b): the page's (P_n)_{<=xi} = P_{n'} cannot hold with the
%   margin's P_n = I*_{n+1}; the reading states (P_n)_{<xi} = P_{n'-1}.
%
% Modern names given and footnoted as ours: skeletal category; direct
% category; semi-simplicial, semi-cubical and globular sets (the presheaf
% categories of the three examples, the third being presheaves on the globe
% category); face poset of a regular CW complex and Björner's CW posets;
% regular polytopes; weak equivalence of higher groupoids for the conditions
% of page 4.
%
% Revised 2026-09-23 (Opus 5.5), after the find-novelty pass:
% — page 11, question b): (P_n)_{<=xi} = P_{n'} replaced by
%   (P_n)_{<xi} = P_{n'-1}, with a footnote on why the page's form fails.
% — page 11, Coxeter footnote: no longer restricted to convex polytopes; the
%   hemispherical disc is the spherical tessellation {2,...,2} and
%   {p,2,...,2} gives further sequences, so the question has a positive
%   answer outside convexity (the edition's remark; Coxeter and McMullen-
%   Schulte cited from memory).
% — pages 2-3: a, l, r (and sigma, mu) marked as the edition's reading; the
%   silent F-dot x I of the amalgamated sum now footnoted against the
%   transcription's B-dot; the struck last factor of F_n said to be struck.
% — page 4: the repeated Z[F_{n-1}] term footnoted.
% — page 5: the typescript names the D_n (not the D-tilde_n).
% — page 10, Ex. 2: the margin's « sections ... Delta_{n-1} x {+-1} -->
%   Delta_{n-1} » footnoted.
% — wording that ordered the leaves in time (« premier état », « réduira »,
%   « essayées d'abord », « prépare la suite », « rédaction au propre ») and
%   one about his intent (« le cas qu'il a en tête ») rephrased to what the
%   pages show.
%
% What the folder announces and never establishes: the proofs of the
% equivalences of pages 9 and 13, both asserted; the sphericity of |I_n^*|
% in the three examples, asserted; the answer to the question of page 11
% (« Ça doit être dans les livres de Coxeter ! »); the canonicity asked in the
% margin of page 11.
\documentclass[11pt,a4paper]{article}
% Le préambule commun à toutes les transcriptions du fonds Grothendieck.
%
% Il tient en une idée : **l'incertitude se compose, elle ne se résout pas.**
% Une transcription qui lisse un mot douteux a détruit la seule chose pour
% laquelle on la faisait — un lecteur doit pouvoir savoir, à chaque endroit,
% ce qui était sur le feuillet et ce qui a été ajouté. Les six macros qui
% suivent sont donc l'appareil critique, et non de la décoration.
%
% Les mêmes commandes sont reconnues par `scripts/render.mjs`, qui produit la
% vue de lecture du site. Une macro ajoutée ici doit l'être aussi là-bas,
% faute de quoi le rendu échouera bruyamment — ce qui est voulu : un rendu
% silencieusement faux est pire qu'un rendu absent.

\ProvidesPackage{grothendieck}[2026/08/08 Transcriptions du fonds Grothendieck]

\usepackage{amsmath,amssymb,amsthm}
\usepackage{mathrsfs}
\usepackage{stmaryrd}
% Les diagrammes commutatifs sont partout dans ce fonds : c'est la langue
% même de la géométrie algébrique de Grothendieck, et une transcription qui
% les rendrait en prose ne transcrirait rien.
\usepackage{tikz-cd}
% Français par défaut — c'est la langue des feuillets — mais l'anglais est
% chargé aussi : la traduction et le résumé sont des documents anglais, et la
% typographie française (espaces avant les deux-points) y serait une faute.
% Ces documents basculent par \selectlanguage{english} en tête de corps.
\usepackage[english,french]{babel}
\usepackage{fontspec}
\usepackage{geometry}
\geometry{a4paper,margin=25mm,marginparwidth=28mm}
\usepackage{marginnote}
\usepackage{xcolor}
% ulem, not soul. `soul`'s \st and \ul are built on a character-by-character
% scan that XeLaTeX's font machinery breaks: the document fails with an
% undefined control sequence rather than degrading. ulem does the same two jobs
% and is Unicode-engine safe. `normalem` stops it hijacking \emph, which the
% transcriptions use for Grothendieck's own emphasis.
\usepackage[normalem]{ulem}
\usepackage{hyperref}
\hypersetup{colorlinks=true,linkcolor=black,urlcolor=black,pdfborder={0 0 0}}

% --- Métadonnées du lot -----------------------------------------------------
% Reprises telles quelles par le script de rendu, qui les affiche en tête de
% la vue de lecture. Elles fixent ce que le fichier prétend couvrir : sans
% elles, une transcription flotte sans point d'attache dans un fonds de seize
% mille feuillets.

\newcommand{\folder}[1]{\gdef\grothFolder{#1}}
\newcommand{\batch}[1]{\gdef\grothBatch{#1}}
\newcommand{\leaves}[2]{\gdef\grothFirst{#1}\gdef\grothLast{#2}}
\newcommand{\foldertitle}[1]{\gdef\grothFoldertitle{#1}}
\gdef\grothFolder{?}\gdef\grothBatch{?}\gdef\grothFirst{?}\gdef\grothLast{?}\gdef\grothFoldertitle{}

% --- Le résumé --------------------------------------------------------------
%
% Il ouvre la lecture modernisée, sous le titre « Résumé », et il est composé
% en petit. Ce n'est pas un ornement : c'est le seul passage du document qui
% ne soit pas des mathématiques, et un lecteur doit pouvoir le reconnaître
% avant de l'avoir lu — puis le sauter s'il connaît déjà le sujet, ou s'y
% arrêter s'il ne le connaît pas. Le filet qui le suit marque la reprise du
% corps.

\newenvironment{resume}
  {\small\setlength{\parskip}{0.45em}}
  {\par\medskip\hrule\medskip}

% --- Mots-clés ---------------------------------------------------------------
%
% En anglais, en clôture du résumé de la lecture modernisée : le vocabulaire
% moderne sous lequel chercher ce que les feuillets construisent. Le manifeste
% du site (`npm run manifest`) les extrait pour étiqueter la cote entière —
% cette ligne est la source unique des tags, il n'y a pas de fichier à côté.
\newcommand{\keywords}[1]{\par\smallskip\noindent
  {\footnotesize\textsc{Keywords} --- \textit{#1}}\par}

% --- L'appareil critique ----------------------------------------------------

% Le numéro de feuillet, en marge. C'est l'ancre qui permet de régler un
% désaccord sur une formule en regardant *un* feuillet plutôt qu'une chemise
% de six cents. Le site s'en sert aussi pour tourner les pages du fac-similé
% quand on fait défiler la transcription.
\newcommand{\leaf}[1]{%
  \marginnote{\footnotesize\color{gray!70}\textsf{#1}}%
  \label{leaf:#1}\ignorespaces}

% Illisible : on ne devine pas. Un mot inventé qui se lit comme les autres est
% le pire résultat possible.
\newcommand{\ill}{\textcolor{red!60!black}{[\,\ldots\,]}}

% Lecture douteuse : le mot est donné, et signalé comme incertain.
\newcommand{\uncertain}[1]{\uline{#1}}

% Ajout de l'éditeur — une lettre manquante, un mot sous-entendu.
\newcommand{\add}[1]{\textcolor{blue!50!black}{[#1]}}

% Biffé par Grothendieck lui-même. Ce qu'il a rayé fait partie du document :
% une rature dit souvent où la pensée a changé de direction.
\newcommand{\struck}[1]{\sout{#1}}

% Note du transcripteur, sur l'état matériel du feuillet.
\newcommand{\note}[1]{\par\smallskip{\small\color{gray!60!black}\textit{[#1]}}\par\smallskip}

% Note marginale de Grothendieck, distincte du corps du texte.
\newcommand{\marginal}[1]{\par\smallskip\noindent\hspace{1em}%
  {\small\color{gray!60!black}#1}\par\smallskip}

% --- Titre ------------------------------------------------------------------

\AtBeginDocument{%
  \begin{center}
    {\large\bfseries Fonds Alexandre Grothendieck --- cote n\textsuperscript{o}~\grothFolder}\\[2pt]
    {\itshape\grothFoldertitle}\\[4pt]
    {\small feuillets \grothFirst--\grothLast\ (lot \grothBatch)}\\[2pt]
    {\footnotesize\color{gray!60!black}Transcription établie d'après le fac-similé numérisé
     par l'Université de Montpellier.\\
     Les lectures incertaines sont soulignées, les illisibles notées [\,\ldots\,],
     les ajouts entre crochets.}
  \end{center}
  \vspace{1em}}

\endinput


\folder{104}
\pages{1}{16}
\dating{s.d.}
\watermark{Édition de démonstration\\interprétation personnelle de l'œuvre}
\foldertitle{[Champs (stacks) 1] : copies de notes manuscrites (s.d.) — lecture modernisée du dossier entier}

\begin{document}

\section*{Résumé}

\begin{resume}
Pour décrire une forme géométrique par des données finies, on la découpe en
cellules élémentaires --- des triangles et leurs analogues en dimension
supérieure, ou des carrés et des cubes, ou encore des disques --- et l'on note
comment elles sont recollées. Le choix de la cellule de base n'est pas
anodin : les triangles donnent les ensembles simpliciaux de la topologie
algébrique, les cubes les ensembles cubiques, et une troisième forme, le
disque découpé en deux hémisphères à chaque dimension, donne les ensembles
« globulaires » dont Grothendieck se sert dans \emph{Pursuing Stacks} pour
parler de groupoïdes de dimension supérieure. Chaque choix fournit une « catégorie de
modèles ».

Ces feuillets posent la question inverse : qu'est-ce qui, dans une catégorie
de modèles, est vraiment nécessaire ? Leur réponse est qu'une telle catégorie
--- sous deux hypothèses simples, que ses flèches soient des inclusions et
qu'un modèle n'ait pas de symétrie propre --- est entièrement déterminée par
des ensembles ordonnés : pour chaque modèle, l'ensemble ordonné de ses faces,
chaque face étant marquée par le modèle auquel elle ressemble. Le triangle de
dimension $n$ a $2^{n+1} - 1$ faces non vides, le cube $3^{n}$, le disque
hémisphérique $2n + 1$ ; et dans les trois cas l'ensemble des faces propres,
réalisé géométriquement, est une sphère.

La seconde moitié applique la même idée aux objets construits en recollant
des modèles. Un tel objet, quand chaque cellule y est plongée et que deux
cellules se rencontrent le long de cellules, se décrit lui aussi par
l'ensemble ordonné de ses cellules, muni d'un « épinglage » qui identifie
chaque cellule et ses faces à un modèle. Les sommes et les recollements le long
de sous-objets se lisent sur ces ensembles ordonnés.

Au passage, une question que le dossier ne tranche pas : existe-t-il
d'autres suites de formes aussi régulières que le simplexe, le cube et le
disque, dont toutes les faces sont encore de la suite ? « Ça doit être dans
les livres de Coxeter ! »

\keywords{direct category, skeletal category, semi-simplicial set, cubical set, globular set, face poset, regular CW complex, CW poset, presheaf category}
\end{resume}

\section*{Le fil du dossier, et les conventions}

\begin{enumerate}
\item[1.] \emph{Feuillets d'essai} (pages 1 à 5) : cinq feuillets en travers,
où apparaissent les mêmes objets de loin --- des systèmes globulaires, des
conditions d'équivalence, et des données de la construction sous une forme
moins rédigée.
\item[2.] \emph{D'une catégorie de modèles à des ensembles ordonnés} (pages 6 à
9) : les données $(\mathbb{N}, I_{n}, \tau_{n}, \varphi_{n,n'})$ et la
reconstruction de $M$.
\item[3.] \emph{Réciproque et exemples} (pages 10 et 11) : simplexe, cube,
disque hémisphérique, et la question des suites régulières.
\item[4.] \emph{Objets admissibles et ensembles ordonnés $M$-épinglés} (pages
12 à 16).
\end{enumerate}

Les pages 6 à 16 sont une rédaction suivie, à l'encre, sous sa pagination 1 à 6,
portée sur les rectos seulement, les versos continuant le texte ; elle est
continue. Rien ne date les feuillets des pages 1 à 5 par rapport à elles.

\emph{Conventions.} $M$ est la catégorie des modèles, $\widehat{M}$ celle des
préfaisceaux d'ensembles sur $M$, dans laquelle $M$ se plonge par
Yoneda.\footnote{Il écrit $\widehat{M}$ à la page 12 et $\widehat{C}$ aux pages
14 et 16, pour le même objet ; on écrit partout $\widehat{M}$.} On note
$\mathbb{N}$ l'ensemble ordonné des classes d'isomorphisme d'objets de $M$, à
ne pas confondre avec l'ensemble $\mathbf{N}$ des entiers naturels.\footnote{Il
distingue les deux par la graisse du trait : un $N$ barré pour les classes, un
$N$ ordinaire pour les entiers, et c'est tout le sens de sa remarque de la page
10, « dans les cas que j'ai en tête, $\mathbb{N} = \mathbf{N}$ ».} Pour
$n \in \mathbb{N}$, $D_{n}$ est un modèle de classe $n$, $I_{n}$ l'ensemble
ordonné de ses sous-objets, $d_{n}$ son plus grand élément,
$\tau_{n} : I_{n} \to \mathbb{N}$ l'application type,
$I_{n,n'} = \tau_{n}^{-1}(n')$, et $u^{n}_{\xi} : I_{n'} \to I_{n}$
l'application admissible attachée à $\xi \in I_{n,n'}$. Pour un ensemble ordonné
$K$ et $\xi \in K$, $K_{\leq \xi} = \{x \in K \mid x \leq \xi\}$.

\emph{Ce que le dossier annonce et n'établit pas} : les démonstrations des deux
équivalences de catégories (pages 9 et 13), toutes deux affirmées ; la
sphéricité de $|I_{n}^{*}|$ dans les trois exemples ; la réponse à la question
de la page 11 et le caractère canonique demandé dans sa marge.

\pagerange{1}{5}
\section*{Feuillets d'essai (pages 1 à 5)}

Les cinq premiers feuillets sont écrits sur du papier de rencontre --- le dos
d'une enveloppe, le dos de pages dactylographiées --- et couverts de figures de
cellules : segments subdivisés, amandes hachurées accolées, carrés divisés, un
huit. Ils ne sont pas rédigés ; on en relève ce qui touche à la suite ou éclaire
le contexte.

\emph{Page 1.} Une chaîne de catégories,
\[
(\text{ensembles simpliciaux}) \supset (\mathrm{Cat}) \longleftarrow
(\text{préordonnés}) \supset (\text{ordonnés}) \supset (\text{Maquettes}),
\]
le dernier mot remplaçant « Triangulations », biffé, et une famille $F_{x} \to
F_{y}$ indexée par $x \leq y$.\footnote{Les inclusions sont tracées d'une
seule courbe, et seule la flèche qui entre dans $(\mathrm{Cat})$ porte une
pointe. Il écrit « Ss sets », « ensembles semi-simpliciaux », ancien nom des
ensembles simpliciaux, avec dégénérescences. La première est le foncteur nerf, pleinement fidèle ; la seconde
regarde un préordre comme une catégorie.}

\emph{Pages 2 et 3 : systèmes globulaires.} Un système
$F_{0} \leftleftarrows F_{1} \leftleftarrows F_{2}$ muni de sources et buts
$s_{i}, t_{i}$, d'unités $\sigma_{i}$, de compositions $\mu_{i}^{j}$ et de trois
lettres $a$, $\ell$, $r$, qu'on peut lire comme les contraintes d'associativité
et d'unité à gauche et à droite d'une structure de dimension
2.\footnote{Cette lecture est la nôtre : la page porte les trois lettres l'une
sous l'autre, sans rien qui dise ce qu'elles désignent. Les rôles d'unités et
de compositions donnés aux $\sigma_{i}$ et aux $\mu_{i}^{j}$ sont de même
lus sur la notation et le sens des flèches, non sur un texte.} Puis des cellules
construites sur une suite de catégories $C_{1}, C_{2}, \ldots$ :
\[
F_{0} = \mathrm{Ob}\, C_{1}, \quad
F_{1} = \mathrm{Fl}\, C_{1} \times \mathrm{Ob}\, C_{2}, \quad
\ldots, \quad
F_{n} = \mathrm{Fl}\, C_{1} \times \cdots \times \mathrm{Fl}\, C_{n}
\times \mathrm{Ob}\, C_{n+1},
\]
la page 2 portant le cas $C_{1} = \Delta_{n}$, $C_{2} = \Delta_{n'}$.\footnote{Les
feuillets s'écartent trois fois de la forme que la série impose, et on la
rétablit : $F_{0} = \mathrm{Ob}\,\Delta_{n'}$ à la page 2, où la suite
demande $\Delta_{n}$ ; $\mathrm{Fl}\, C_{2}$ en tête de la troisième ligne
de la page 3, où elle demande $C_{1}$ ; et, à la ligne de $F_{n}$, un dernier
facteur $\times\, \mathrm{Ob}\, C_{4}$ biffé, la ligne s'arrêtant à
$\mathrm{Fl}\, C_{n}$, où la suite demande $\times\, \mathrm{Ob}\,
C_{n+1}$. La page 3 porte aussi des sous-ensembles $O_{n} \subset F_{n}$, un
cylindre $(F \times I) \sqcup_{\dot{F} \times I} \dot{F}$ et un produit de
joints $(O_{1}, F_{1}) * (O_{2}, F_{2}) * \cdots$, sans texte. Dans le
cylindre, la transcription lit l'indice $\dot{B} \times I$, avec doute ; aucun
$B$ n'est défini sur la page, et la somme amalgamée n'a de sens que le long de
$\dot{F} \times I \subset F \times I$ et de la projection $\dot{F} \times I
\to \dot{F}$, d'où notre $\dot{F}$.}

\emph{Page 4 : quand un morphisme est-il une équivalence.} La suite
$\mathbb{Z}[F_{n+1}] \to \mathbb{Z}[F_{n}] \xrightarrow{t_{n} - s_{n}}
\mathbb{Z}[F_{n-1}]$ est un complexe\footnote{La page écrit à gauche un
quatrième terme, $\mathbb{Z}[F_{n-1}] \longleftarrow
\mathbb{Z}[F_{n-1}]$, avec deux fois le même indice ; le terme suivant de la
suite serait $\mathbb{Z}[F_{n-2}]$. On ne garde que les trois termes que
l'identité ci-dessous concerne.}, parce que
$t_{n-1}(t_{n}y - s_{n}y) - s_{n-1}(t_{n}y - s_{n}y) = 0$ par les identités
globulaires $s_{n-1}s_{n} = s_{n-1}t_{n}$ et $t_{n-1}s_{n} = t_{n-1}t_{n}$.
Puis, pour $f : F \to F'$, des conditions échelonnées : a) tout objet de
$F'$ est relié par une flèche à l'image d'un objet de $F$ ; b) deux objets
dont les images sont reliées le sont --- « $\pi_{0}(f)$ inj. » ; b') et on peut
choisir la flèche relevée à une 2-flèche près ; c) et c') les mêmes énoncés une
dimension plus haut --- « $\pi_{1}(f)$ inj. ». Une colonne marginale en donne la
forme générale :
\begin{enumerate}
\item[$(C_{i})$] pour tout $w' \in F'_{i}$, il existe $w \in F_{i}$ tel que
$f(w) \sim w'$, et, si $i \geq 1$, on peut se donner d'avance la source et le
but de $w$, sous la seule condition que leurs images soient ceux de
$w'$.\footnote{C'est le feuillet le moins lisible du dossier, et la colonne
marginale ne se lit que par fragments. La dernière ligne de c') écrit
« $w \sim f(w)$ », où le sens demande $w' \sim f(w)$. La condition $(C_{i})$ est
la forme, à bord prescrit, de la définition des équivalences faibles de
groupoïdes de dimension supérieure ; le nom n'est pas sur la page.}
\end{enumerate}

\emph{Page 5 : les données de la construction, en abrégé.} Au dos d'une page
de son propre tapuscrit anglais sur les « Gr-stacks », des données que les
pages 6 à 9 rédigent, l'ensemble ordonné des classes s'appelant ici $I$ au lieu
de $\mathbb{N}$ :
une famille d'ensembles ordonnés $n \mapsto \widetilde{D}_{n}$, une application
strictement croissante $\tau_{n} : \widetilde{D}_{n} \to I$ d'image $I_{\leq n}$,
et le dénombrement des faces du cube
\[
1 + n \cdot 2 + \frac{n(n-1)}{2}\, 2^{2} + \cdots + 2^{n} = (1 + 2)^{n}
\]
que la marge de la page 10 écrit $\mathrm{Card}\, I_{n} = 3^{n}$.\footnote{La page note
aussi « $\pi_{0}(F_{*}, x)$ = ens.\ des classes de lacets en $x$ », avec un lacet
dessiné ; des classes de lacets en un point forment un $\pi_{1}$. Le mot
« classes » est lu avec doute. Le tapuscrit anglais du verso pose
$\underline{B} = (\text{Gr-stacks})$ et appelle $D_{n}$ (souligné à la
machine) les objets de base, dont les $\widetilde{D}_{n}$ de la page
manuscrite sont les ensembles de sous-objets ; le terme
est celui de \emph{Pursuing Stacks}, rédigé en anglais en 1983, ce qui est un
rapprochement et non une datation.}

\pagerange{6}{9}
\section*{D'une catégorie de modèles à des ensembles ordonnés (pages 6 à 9)}

Soit $M$ une catégorie telle que
\begin{enumerate}
\item[(i)] toute flèche de $M$ est un monomorphisme ;
\item[(ii)] tout endomorphisme d'un objet de $M$ est l'identité.
\end{enumerate}
Exemple : la catégorie définie par un ensemble ordonné.\footnote{« Tout
endomorphisme » remplace « tout automorphisme », biffé. La correction
importe : avec (i) seulement et des automorphismes triviaux, un endomorphisme
non trivial resterait possible dans une catégorie infinie. Une catégorie qui
satisfait (i) et (ii) est une \emph{catégorie directe} au sens de Reedy, dont
toutes les flèches sont monomorphes, dès que l'ordre de $\mathbb{N}$ défini
plus bas est bien fondé ; le nom n'est pas sur la page.}

\emph{Sous-objets.} Pour $D \in \mathrm{Ob}\, M$, soit $\widetilde{D}$
l'ensemble des sous-objets de $D$, ordonné par inclusion, et identifions $D$ au
plus grand élément. Pour $D', D$ objets de $M$, l'application
\[
(*) \qquad \mathrm{Hom}_{M}(D', D) \longrightarrow \widetilde{D},
\qquad u \longmapsto u(D')
\]
est injective : si $u$ et $v$ définissent le même sous-objet, $v = u\theta$
avec $\theta$ automorphisme de $D'$, et $\theta = \mathrm{id}$ par (ii). Son
image est l'ensemble des sous-objets de $D$ isomorphes à $D'$.

\emph{Les classes.} L'ensemble $\mathbb{N}$ des classes d'isomorphisme
d'objets de $M$, préordonné par l'existence d'une flèche, est ordonné : deux
flèches $D' \to D \to D'$ composent en un endomorphisme, donc en
l'identité, et $D \simeq D'$. Pour $n \in \mathbb{N}$, un objet $D_{n}$ de
classe $n$ est déterminé à isomorphisme unique près, et $I_{n} =
\widetilde{D}_{n}$ est donc défini à isomorphisme unique près.

\emph{Le type} (page 7). L'application $\tau_{n} : I_{n} \to \mathbb{N}$, qui
à un sous-objet associe sa classe, est strictement croissante, d'image
$\mathbb{N}_{\leq n}$, et
\[
\tau_{n}^{-1}(n) = \{d_{n}\}.
\]
En effet, si $x \subset y$ sont deux sous-objets de même type, la flèche
$x \to y$ admet une rétraction et le composé dans l'autre sens est un
endomorphisme de $y$, donc l'identité : $x = y$.\footnote{La page énonce la
stricte croissance et $\tau_{n}^{-1}(n) = \{d_{n}\}$ « par (ii) » ; l'argument
est le nôtre. « Rappelle » et « déjà » sont lus sur le sens.}

\emph{Les applications admissibles.} Pour $n' \leq n$, $(*)$ donne une bijection
$\mathrm{Hom}(D_{n'}, D_{n}) \simeq I_{n,n'} = \tau_{n}^{-1}(n')$, et un
monomorphisme $u : D_{n'} \to D_{n}$ induit par image directe une application
$I_{n'} \to I_{n}$. D'où
\[
\varphi_{n,n'} : I_{n,n'} \longrightarrow \mathrm{Hom}(I_{n'}, I_{n}),
\qquad \xi \longmapsto u^{n}_{\xi},
\]
avec, pour tout $\xi \in I_{n,n'}$ (page 8),
\[
u^{n}_{\xi}(d_{n'}) = \xi, \qquad
u^{n}_{\xi} : I_{n'} \xrightarrow{\;\sim\;} (I_{n})_{\leq \xi}, \qquad
u^{n}_{d_{n}} = \mathrm{id}_{I_{n}},
\]
de sorte que $\varphi_{n,n'}$ est injective ; et, pour
$n'' \leq n' \leq n$, $\xi \in I_{n,n'}$, $y \in I_{n',n''}$, la \emph{loi de
composition}
\[
u^{n}_{\xi} \circ u^{n'}_{y} = u^{n}_{\zeta}, \qquad \zeta = u^{n}_{\xi}(y),
\]
qui exprime la stabilité des applications admissibles par composition.

\emph{La catégorie $M_{0}$.} Elle a pour objets les éléments de $\mathbb{N}$,
pour flèches $n' \to n$ les applications admissibles $I_{n'} \to I_{n}$ si
$n' \leq n$ et aucune sinon, la composition étant celle des applications. On a
un foncteur $M \to M_{0}$, $D \mapsto$ sa classe, $u \mapsto$ l'image directe,
et c'est une équivalence de catégories ; c'est un isomorphisme si $M$ est
squelettique.\footnote{Il écrit « $M$ est une catégorie \emph{rigide} (deux
objets isom.\ sont identiques) » : la définition entre parenthèses est celle
qu'on appelle aujourd'hui \emph{squelettique}, « rigide » désignant d'ordinaire
l'absence d'automorphismes non triviaux, que (ii) assure déjà. La page ne
démontre pas l'équivalence ; elle est fidèle par $(*)$, pleine par définition
de $M_{0}$, et essentiellement surjective par construction.}

Supposant désormais $M$ squelettique, on l'identifie à $M_{0}$, catégorie
définie par les données suivantes :
\begin{itemize}
\item $\mathbb{N}$, ensemble ordonné ;
\item $(I_{n})_{n \in \mathbb{N}}$, ensembles ordonnés ayant chacun un plus grand
élément $d_{n}$ ;
\item $(\tau_{n} : I_{n} \to \mathbb{N})$, strictement croissantes,
$\tau_{n}(I_{n}) = \mathbb{N}_{\leq n}$, $\tau_{n}^{-1}(n) = \{d_{n}\}$ ;
\item $(\varphi_{n,n'})_{n' \leq n}$, $\xi \mapsto u^{n}_{\xi}$, avec
$u^{n}_{\xi} : I_{n'} \simeq (I_{n})_{\leq \xi}$ et la loi de composition.
\end{itemize}

\pagerange{10}{11}
\section*{Réciproque, exemples, et une question (pages 10 et 11)}

\emph{Réciproque.} De telles données définissent une catégorie $M_{0}$
satisfaisant à (i) et (ii), dont les ensembles de sous-objets sont les
$I_{n}$.\footnote{Deux propriétés qu'on vient d'établir pour une catégorie
donnée ne figurent pas dans la liste des données : $u^{n}_{d_{n}} =
\mathrm{id}$ et $\tau_{n} \circ u^{n}_{\xi} = \tau_{n'}$. Elles s'en
déduisent. La loi de composition donne $u_{d_{n}} \circ u_{d_{n}} = u_{d_{n}}$,
et $u_{d_{n}}$ étant bijective, c'est l'identité. Elle impose aussi que
$u^{n}_{\zeta}$ ait pour source $I_{n''}$, c'est-à-dire
$\tau_{n}(u^{n}_{\xi}(y)) = \tau_{n'}(y)$. Les flèches de $M_{0}$ sont
injectives, d'où (i), et $\mathrm{End}(n)$ s'identifie à
$\tau_{n}^{-1}(n) = \{d_{n}\}$, d'où (ii). Enfin $\xi \leq \eta$ dans $I_{n}$
si et seulement si $u_{\xi}$ se factorise par $u_{\eta}$, parce que
$(I_{n})_{\leq \eta}$ est l'image de $u_{\eta}$ : les sous-objets de $n$ dans
$M_{0}$ forment bien $I_{n}$. Ces vérifications sont les nôtres.} « Dans les cas
que j'ai en tête, $\mathbb{N} = \mathbf{N}$. »

\textbf{Exemple 1.} La catégorie des ensembles finis totalement ordonnés
$\Delta_{n} = \{0, \ldots, n\}$ et des applications strictement croissantes.
Les sous-objets sont les parties non vides :
\[
I_{n} = \mathfrak{P}^{*}(\Delta_{n}), \qquad
\mathrm{Card}\, I_{n} = 2^{n+1} - 1 .
\]

\textbf{Exemple 2.} La catégorie des cubes combinatoires épinglés --- par un
ordre total sur l'ensemble des paires d'hyperfaces opposées et le choix d'une
face dans chaque paire --- et des inclusions de faces compatibles aux
épinglages, l'épinglage d'un cube induisant celui de chacune de ses faces.
Une face de $[0,1]^{n}$ fixe certaines coordonnées à $0$ ou $1$ et laisse les
autres libres :
\[
\mathrm{Card}\, I_{n} = \sum_{k} \binom{n}{k} 2^{k} = 3^{n}.
\]
Autrement dit, $I_{n}$ est l'ensemble des sections partielles de la projection
$\Delta_{n-1} \times \{\pm 1\} \to \Delta_{n-1}$ : une face choisit, pour
certaines des $n$ paires d'hyperfaces, l'une des deux.\footnote{La note
marginale de l'Ex. 2 commence par « $I_{n}$ = ens.\ des sections », le mot
« sections » lu avec doute, et porte $\Delta_{n-1} \times \{\pm 1\} \to
\Delta_{n-1}$ ; le reste n'est lu que par fragments. Le mot « partielles » et
l'ordre (une face est inférieure à une autre si sa section prolonge
l'autre) sont les nôtres.}

\textbf{Exemple 3.} La catégorie des disques hémisphériques $D_{n}$, dont les
cellules sont deux hémisphères en chaque dimension $k < n$ et la cellule
intérieure :
\[
I_{n} = \bigl(\Delta_{n-1} \times \{\pm 1\}\bigr) \sqcup \{d_{n}\},
\qquad \mathrm{Card}\, I_{n} = 2n + 1,
\]
ordonné par $(k, \varepsilon) < (k', \varepsilon')$ si et seulement si
$k < k'$, et $d_{n}$ plus grand élément.\footnote{Le corps de la page écrit
$\Delta_{n} \times \{\pm 1\}$, qui aurait $2n + 3$ éléments ; la marge écrit
« Card $I_{n} = 2n+1$ », et c'est ce que donne la géométrie : un disque de
dimension $n$ a pour bord une sphère $S^{n-1}$ découpée en deux hémisphères en
chacune des dimensions $0, \ldots, n-1$. On retient $\Delta_{n-1}$. La page
munit $\{\pm 1\}$ de « l'ordre chaotique » ; le produit par un préordre
chaotique serait un préordre et non un ordre, deux hémisphères de même
dimension y étant chacun inférieur à l'autre. L'ordre voulu, celui des cellules
du disque, est celui qu'on écrit.}

\emph{Ce que sont ces trois catégories.} Les préfaisceaux sur la première sont
les ensembles semi-simpliciaux au sens actuel du mot, c'est-à-dire sans
dégénérescences, puisque toutes les flèches sont monomorphes ; sur la deuxième, des ensembles semi-cubiques ; sur
la troisième, les ensembles globulaires. En effet $\mathrm{Hom}(D_{n'}, D_{n})
= \tau_{n}^{-1}(n')$ a deux éléments pour $n' < n$, les deux hémisphères de
dimension $n'$, et c'est la catégorie des globes, avec sa source et son
but.\footnote{L'identification est la nôtre ; la page ne nomme aucune des
trois. Le troisième exemple est la forme des cellules de \emph{Pursuing
Stacks}. Le mot « hémisphères » est lu avec doute.}

\emph{Sphéricité.} Dans les trois exemples, posant $I_{n}^{*} = I_{n} -
\{d_{n}\}$ :
\begin{enumerate}
\item[a)] pour $n \geq 1$, la réalisation géométrique $|I_{n}^{*}|$ est une
sphère de dimension $n - 1$, et $|I_{n}|$ un disque de dimension $n$ ;
\item[b)] les $I_{n}^{*}$ sont \emph{réguliers} : leur groupe d'automorphismes
est transitif sur les repères, c'est-à-dire sur les chaînes maximales
$\xi_{0} < \xi_{1} < \cdots < \xi_{n-1}$.
\end{enumerate}
Le groupe est $\mathfrak{S}_{n+1}$ pour le simplexe, le groupe hyperoctaédral
pour le cube, et $\{\pm 1\}^{n}$ pour le disque hémisphérique, qui y opère
simplement transitivement.\footnote{Les trois groupes sont les nôtres ; la page
énonce a) et b) sans démonstration. La réalisation d'un ensemble ordonné est
celle de son complexe des chaînes. Un ensemble ordonné dont chaque
$K_{< \xi}$ se réalise en une sphère est ce qu'on appelle aujourd'hui un
\emph{CW-poset} : c'est exactement l'ensemble ordonné des cellules d'un
complexe CW régulier (A.~Björner, 1984). Les trois exemples en sont.}

\textbf{Question} (page 11). Existe-t-il d'autres suites $(P_{n})_{n \geq 1}$
d'ensembles ordonnés telles que
\begin{enumerate}
\item[a)] $P_{n}$ soit un polyèdre sphérique régulier ;
\item[b)] pour tout $\xi \in P_{n}$, $(P_{n})_{< \xi} \simeq P_{n'-1}$ pour un
$n' \leq n$ convenable --- toute face fermée de $P_{n}$ est un disque dont le
bord est un terme de la suite ?
\end{enumerate}
« Ça doit être dans les livres de Coxeter ! » On aimerait surtout pouvoir
construire des $\varphi_{n,n'}$ satisfaisant la loi de
composition.\footnote{En marge : « b) est-il canonique ?! », et
$P_{n} = I_{n+1}^{*}$, ce qui fixe le décalage d'indice entre la question et les
exemples. La page écrit b) sous la forme $(P_{n})_{\leq \xi} \simeq P_{n'}$,
qui ne peut pas tenir : $(P_{n})_{\leq \xi}$ a un plus grand élément, $\xi$,
et sa réalisation est un cône, donc contractile, alors que $|P_{n'}|$ est une
sphère par a). Avec $P_{n} = I_{n+1}^{*}$, si $\xi$ est de type $n'$,
$(P_{n})_{\leq \xi} = (I_{n+1})_{\leq \xi} \simeq I_{n'}$ par
l'application admissible de la page 8, et son bord est $I_{n'}^{*} =
P_{n'-1}$ ; c'est ce qu'on écrit, en convenant que $P_{0}$ est la sphère de
dimension $0$ (deux points) et $P_{-1}$ l'ensemble vide. La forme de la page
est juste pour les disques $I_{n}$ et non pour les sphères $P_{n}$.
Parmi les polytopes réguliers convexes, la condition b) --- toute face
est encore de la suite --- ne laisse passer indéfiniment que les simplexes et
les cubes : les faces d'un octaèdre sont des triangles, et une suite qui
contiendrait l'octaèdre devrait se poursuivre par un polytope régulier de
dimension 4 à cellules octaédriques, le 24-cellules, qui n'a pas de successeur.
Mais la question ne demande pas la convexité, et sans elle la réponse est oui.
Le disque hémisphérique est lui-même un pavage régulier de la sphère, de
symbole de Schläfli $\{2, 2, \ldots, 2\}$ : deux hémisphères en chaque
dimension. Et pour tout $p \geq 2$, la suite point, segment, $p$-gone, puis,
en chaque dimension suivante, le disque dont le bord est formé de deux
exemplaires du terme précédent recollés le long de leur bord --- les pavages
$\{p, 2, \ldots, 2\}$ --- est régulière et a toutes ses faces dans la
suite ; $p = 2$ redonne le disque hémisphérique. Coxeter
traite les pavages « dégénérés » de la sphère, dièdres $\{p, 2\}$ et
hosoèdres $\{2, p\}$, dans \emph{Regular Polytopes}, et les symboles de ce
genre relèvent aujourd'hui des polytopes réguliers abstraits (P.~McMullen et
E.~Schulte) : deux références citées de mémoire, sans vérification de page.
Que ces suites portent des $\varphi_{n,n'}$ satisfaisant la loi de composition
n'est pas examiné ici. Cette remarque est la nôtre.}

\pagerange{12}{13}
\section*{Objets admissibles et épinglages (pages 12 et 13)}

Soit $\mathbb{B}_{0}$ la sous-catégorie pleine de $\widehat{M}$ formée des
objets $X$ --- dits \emph{admissibles} --- tels que
\begin{itemize}
\item $X$ est la borne supérieure de ses sous-objets isomorphes à des objets de
$M$ (ses sous-objets \emph{modèles}) ;
\item si $K$ et $L$ sont deux sous-objets modèles de $X$, $K \cap L$ est la
borne supérieure des sous-objets modèles qu'il contient.
\end{itemize}

Soit $K = \widetilde{X}$ l'ensemble ordonné des sous-objets modèles de $X$. Il
porte une structure de \emph{$M$-épinglage} :
\begin{enumerate}
\item[a)] une application strictement croissante
$\tau_{K} : K \to \mathbb{N}$, l'application type ;
\item[b)] pour tout $\xi \in K$, de type $n = \tau_{K}(\xi)$, un isomorphisme
d'ensembles ordonnés
\[
u^{K}_{\xi} : I_{n} \xrightarrow{\;\sim\;} K_{\leq \xi},
\qquad u^{K}_{\xi}(d_{n}) = \xi,
\]
l'application d'épinglage de $K_{\leq \xi}$ ;
\end{enumerate}
qui satisfait
\[
(1) \qquad \tau_{K} \circ u^{K}_{\xi} = \tau_{n},
\]
\[
(2) \qquad u^{K}_{\xi} \circ u^{n}_{y} = u^{K}_{\zeta}
\quad \text{pour } y \in I_{n},\ \zeta = u^{K}_{\xi}(y).
\]
Les sous-objets modèles de $X$ contenus dans $\xi \simeq D_{n}$ sont ceux de
$D_{n}$, par Yoneda et (i), d'où b).\footnote{Il écrit l'application type à
valeurs dans $\mathbf{N}$, les entiers, sur un $\mathbb{N}$ biffé ; c'est le
cas $\mathbb{N} = \mathbf{N}$ que la page 10 dit avoir « en tête ». On garde $\mathbb{N}$, qui est le cas
général. En marge, les applications $I_{n} \to K$ ainsi définies sont dites
« admissibles ».}

\textbf{Énoncé.} La catégorie des objets admissibles de $\widehat{M}$ est
équivalente à celle des ensembles ordonnés $M$-épinglés, dont les morphismes
$f : K \to L$ sont les applications strictement croissantes telles que
\begin{enumerate}
\item[a)] $\tau_{L} \circ f = \tau_{K}$ ;
\item[b)] pour tout $\xi \in K$, avec $y = f(\xi)$ et $n = \tau_{K}(\xi) =
\tau_{L}(y)$, le triangle
\end{enumerate}

\begin{tikzcd}[column sep=small, row sep=small, nodes={font=\scriptsize}]
K_{\leq \xi} \arrow[rr, "f|K_{\leq \xi}"] & & L_{\leq y} \\
& I_n \arrow[ul, "u^K_\xi"] \arrow[ur, "u^L_y"'] &
\end{tikzcd}

commute.\footnote{Le sommet de droite du triangle est écrit comme un $K$ ;
$f$ allant de $K$ dans $L$, c'est $L_{\leq y}$, que la page 14 écrit. Les deux
obliques ne portent sur la feuille qu'un signe d'isomorphie. L'équivalence est
affirmée et non démontrée ; ce qui la rend plausible est que tout élément de
$X(D_{n})$ se factorise par un sous-objet modèle, donc est un monomorphisme, de
sorte qu'un morphisme d'objets admissibles envoie modèles sur modèles, et
qu'un objet admissible est la limite inductive de ses sous-objets modèles
indexée par $K$. Cette esquisse est la nôtre.}

\pagerange{14}{15}
\section*{Parties fermées, et épinglages par les éléments maximaux (pages 14 et 15)}

Un morphisme $f$ comme ci-dessus induit, pour tout $\xi$, un isomorphisme
$K_{\leq \xi} \simeq L_{\leq f(\xi)}$. On s'intéresse surtout au cas où $f$ est
injectif, qui correspond aux monomorphismes de $\widehat{M}$. Alors $f$ est un
isomorphisme de $K$ sur son image $L' = f(K)$, munie de l'ordre induit ; $L'$ est
une partie \emph{fermée} de $L$, c'est-à-dire stable par passage aux éléments
plus petits ; une partie fermée porte un unique $M$-épinglage induit pour
lequel l'inclusion est admissible ; et $f : K \simeq L'$ est compatible aux
épinglages.

\emph{Épinglages par les éléments maximaux} (page 15). Supposons que $K$ n'ait
qu'un nombre fini d'éléments maximaux et que tout élément soit inférieur à
l'un d'eux --- c'est le cas où $X$ est réunion finie de modèles. Un
$M$-épinglage de $K$ est alors connu dès qu'on connaît les
\[
u^{K}_{\alpha} : I_{\tau(\alpha)} \xrightarrow{\;\sim\;} K_{\leq \alpha}
\qquad (\alpha \in \mathrm{Max}\, K),
\]
et une telle famille définit un épinglage si et seulement si, pour
$\alpha, \beta \in \mathrm{Max}\, K$ et $\xi \leq \alpha, \beta$, les deux
isomorphismes $I_{n} \simeq K_{\leq \xi}$ que $u^{K}_{\alpha}$ et
$u^{K}_{\beta}$ induisent --- chacun composé avec l'application admissible de
l'antécédent de $\xi$ --- sont les mêmes.\footnote{Implicitement, les deux
antécédents doivent avoir le même type $n$, sans quoi les deux isomorphismes
n'ont pas la même source ; c'est une partie de la condition. Le dernier mot de
la page, « mêmes », est lu avec doute.} C'est une condition de recollement,
comme pour un atlas.

\pagerange{16}{16}
\section*{Sommes et sommes amalgamées (page 16)}

La catégorie des objets admissibles de $\widehat{M}$ est stable par sommes
quelconques et par sommes amalgamées $X \sqcup_{Z} Y$ le long de
monomorphismes $Z \to X$ et $Z \to Y$, et le foncteur $X \mapsto
\widetilde{X}$ commute à ces opérations. Les objets de type fini --- réunions
finies de modèles --- sont stables par sommes finies et par ces sommes
amalgamées.\footnote{La page dit la commutation « de façon évidente ». Voici
pourquoi elle l'est. Les éléments de $X \sqcup_{Z} Y$ sur $D_{n}$ sont ceux de
$X$ et de $Y$ recollés le long de $Z$, de sorte qu'un sous-objet modèle du
recollement est dans $X$ ou dans $Y$. D'autre part, tout sous-préfaisceau d'un
modèle est réunion de modèles, puisque toutes les flèches de $M$ sont
monomorphes, ce qui assure la seconde condition d'admissibilité. Enfin, un
sous-objet d'un objet admissible est admissible, et $Z$ l'est donc. Ces
vérifications sont les nôtres.}

\textbf{N.B.} Quand les $I_{n}$ sont finis, les objets admissibles de type fini
correspondent aux ensembles ordonnés $M$-épinglés finis.

\end{document}
